\documentclass[letterpaper, 12pt, parskip=full,DIV=10]{scrartcl}

\title{PLSC 40601: Advanced Topics in Causal Inference.}
\subtitle{University of Chicago, Spring 2026.}

\newcommand{\plscassetpath}{assets}
\IfFileExists{assets/template_MOW.sty}{}{%
  \renewcommand{\plscassetpath}{../assets}%
}
\newcommand{\plscbibpath}{assets/PLSC40601}
\IfFileExists{assets/PLSC40601.bib}{}{%
  \renewcommand{\plscbibpath}{../assets/PLSC40601}%
}

\RequirePackage{\plscassetpath/template_MOW}

\begin{document}
\maketitle

\vskip -10 ex

\textbf{Location:} Pick Hall 222\\
\textbf{Course time:} Mon/Wed 16:30--17:50\\
\textbf{Course dates:} March 23, 2026--May 20, 2026

\textbf{Instructor:} Molly Offer-Westort; \href{mailto:mollyow@uchicago.edu}{mollyow@uchicago.edu}\\
\textbf{Office hours:} Book at \url{https://calendar.app.google/DZYbbZHQqCSYAuUr5}\\
\textbf{Office:} Pick Hall 328, 5828 S. University Ave.

\paragraph{Course overview.}
This is a graduate-level course considering modern advances in causal inference and experimental design. 
In particular, we will consider how machine learning methods can be leveraged to address causal questions. 
We will read a selection of papers introducing and implementing techniques and research designs, with applications to the social and health sciences and public policy. 
We will discuss what these new methods are able to offer, and where they may have limitations. 
The course will be oriented around class discussion and student presentations on the readings. 
An introductory course in probability and statistics is required; this prerequisite can be met by courses in statistics, biostatistics, economics, political science, sociology, or related fields. 
Coursework in causal inference is recommended but not required; additional reading references will be provided for students who have not had prior exposure to causal inference methodology. 

\paragraph{Course logistics.}
A public repository is used to distribute the syllabus, slide decks, public code, and other public course materials. 
Required readings are distributed privately and are intentionally not stored in this repository.

\paragraph{AI policy.}
The public AI policy for the course is posted in the course repository as \texttt{ai-policy.md}.
Students remain responsible for the design, interpretation, and verification of all submitted work. 
Individual assignments may impose additional verification requirements. 

\paragraph{Reference readings.}
This class will not follow a textbook, but some of the methods and approaches may be new to students. For reference for machine learning methods, I recommend:
\begin{itemize}
\item \bibentry{hastie2009elements} \vskip 2 ex

A free pdf version of the book is available at the publisher's website: \url{https://link.springer.com/book/10.1007/978-0-387-84858-7}
\end{itemize}


Other machine learning references are:
\begin{itemize}
\item \bibentry{murphy2012machine}
\item \bibentry{bishop2006pattern}
\end{itemize}

Some older reference readings on core techniques in machine learning that are still relevant today--I really enjoy all of these papers and recommend them for students who are interested in the history and development of machine learning methods:
\begin{itemize}
  \item \bibentry{breiman1996bagging}
  \item \bibentry{efron1979bootstrap}
  \item \bibentry{breiman2001random}
\end{itemize}

A causal/{ML} textbook:
\begin{itemize}
\item \bibentry{chernozhukov2024applied}

\url{https://arxiv.org/abs/2403.02467}
\end{itemize}

Many of the papers we read in this course will use potential outcomes notation. Students who have not had recent coursework on causal inference should read the first few chapters of either:
\begin{itemize}
\item
\bibentry{hernan2010causal}

A free pdf version of the book is available at the author's website: \url{https://www.hsph.harvard.edu/miguel-hernan/
causal-inference-book/}
\item 
Or\\
\bibentry{imbens2015causal}
\item Also relevant:\\
\bibentry{ding2023first}

\url{https://arxiv.org/abs/2305.18793}
\end{itemize}

As well, the materials for PLSC 30600: Causal Inference are available at \url{https://github.com/UChicago-pol-methods/PLSC30600}. 

Much of the statistical theory referenced in readings is covered in this book: 
\begin{itemize}
\item \bibentry{lehmann1999elements}
\end{itemize}

Also useful for econometrics references:
\begin{itemize}
\item \bibentry{angrist2009}
\end{itemize}


\paragraph{Presentations.}
Each student will be responsible for presenting and leading discussion on one of the required papers, beginning the second week of the quarter. 
Discussion leaders will provide a 15-20 minute total high-level overview of the paper, and will develop a list of discussion questions. 
For classes with multiple leaders, co-discussion leaders may divide papers among themselves or may collaborate. 
Discussion questions should be submitted to the instructor by email the evening before class. 

\paragraph{Referee report.}
Select one of the papers on the syllabus, including from the additional references. 
Write a referee report as if you were requested to review the paper for publication at a top journal, summarizing the paper's main contributions, areas that could be strengthened, and additional extensions or analyses that you would like to see. 
These are already published papers, so you are unlikely to find significant errors or typos (although these may be present!), but you can evaluate the paper's contribution and novelty; you can imagine that you are a reviewer for the journal where the paper was published. 
In your report, you should also evaluate the paper's clarity and accessibility---it is acceptable to say that there are things that you don't understand. 
Submit the referee report on Canvas \textbf{the night before the paper is discussed in class}. 
The paper selected for the referee report cannot be from the same class as the class presentation. 
References on how to write a referee report:
\begin{itemize}
\item \bibentry{berk2016preparing}
 \item \url{https://chrisblattman.com/blog/2012/01/18/how-to-referee-an-academic-paper/}
 \item \bibentry{thomson2001guide} Chapter 4: \emph{Writing Referee Reports}. 
 \end{itemize}

\paragraph{Final project.}
The final project is an original tutorial in which you explore a topic related to developments in machine learning for causal inference both analytically and computationally. 
Briefly discuss the foundations for this work in statistics or related fields. 
Demonstrate the method with a computational example (i.e. write reproducible code). 
The deliverable is a script or notebook and related data that can be run all the way through after setting the working directory. 
The reader of your tutorial should come away with a good understanding of the strengths and drawbacks of the method / theme as well as some ideas about how to implement.

You will be graded not only on the accuracy of your overview, but also on the clarity of your exposition, which includes notation, typesetting (equation alignment, exceeding margins), grammar, figures, tables, etc. 
All code should be reproducible--please provide any supporting files necessary with your submission. 
The more difficulty I have in reading your work--both text and code--the fewer points you will receive. 
Follow a coding style guide  (e.g. \url{https://style.tidyverse.org/} and \url{https://peps.python.org/pep-0008/}). 
I hope to see polished work in the final product.

Timeline:
\begin{enumerate}
\item Tutorial topics must be emailed to the instructor by the Wednesday of Week 4.  
\item Meet with the instructor to present work-in-progress before the end of Week 7. Send the current state of your work as well as a short email describing your current challenges and areas you would like feedback on 24 hours before the meeting. 
\item Submit the assignment by 11:59 PM on the Friday of finals week.
\end{enumerate}

\begin{itemize}[topsep=0pt,itemsep=-1ex,partopsep=1ex,parsep=1ex]
\item Leading class discussions: 25\% 
\item Referee report: 25\%
\item Final paper: 40\%
\item Participation: 10\%
\end{itemize}

\paragraph{Accommodations:} 
Please reach out to me directly if you would like to request accommodations for the course to better facilitate your learning. 
Student Disability Services (\url{https://disabilities.uchicago.edu/}) is also available to provide you resources and support, and may provide approval for specific academic accommodations. 
Informing me in a timely manner will help me to ensure accommodations are met and I am able to implement an appropriate assessment of your learning. 

\section*{Course outline.}

\subsection*{Week 1. Causal inference, identification, course orientation.}

\subsubsection*{Monday. Lecture: discussion and workflow.}
\begin{itemize}
\item \bibentry{murphy1997read}
\end{itemize}

\subsubsection*{Wednesday. Lecture: course scope.}
\begin{itemize}
\item \bibentry{athey2019machine}
\item \bibentry{breiman2001statistical}
(comments and rejoinder optional)\vskip 2 ex
\item \textit{Textbook reference:} \bibentry{hernan2010causal} A definition of causal effect, 1. 
\end{itemize}

\subsection*{Week 2. Adjustment under selection on observables.}


\subsubsection*{Monday.}
\begin{itemize}
\item \bibentry{bloniarz2016lasso}\vskip 2 ex
\item \textit{Optional/reference:} \bibentry{belloni2014inference}
\item \textit{Optional/reference:} \bibentry{urminsky2016using}
\item \textit{Optional/reference:} \bibentry{stone1974cross}
\vskip 2 ex
\item \textit{Textbook reference:} \bibentry{hastie2009elements} Shrinkage, 3.4; Sample error and cross-validation, 7.2-7.5, 7.10. 
\item \textit{Textbook reference:} \bibentry{hernan2010causal} Randomized experiments 2.1-2.2; Why model, 11.1-11.5; The curse of dimensionality 10.5. 
\end{itemize}

\subsubsection*{Wednesday.}
\begin{itemize}
\item \bibentry{athey2018approximate}\vskip 2 ex
\item \textit{Textbook reference:} \bibentry{hastie2009elements} Bootstrap methods 7.11; Bagging 8.7.
\item \textit{Textbook reference:} \bibentry{hernan2010causal} Observational studies 3.1-3.5; IP weighting and marginal structural models 12.1-12.4.
\end{itemize}

\subsection*{Week 3. Heterogeneous treatment effects.}
\begin{itemize}
\item Trees, forests, and meta-learners for treatment-effect heterogeneity.
\item Distinguishing between estimating a CATE surface, ranking units, and supporting an actionable policy.
\item Interpreting heterogeneity estimates as inputs to later decision-making.
\end{itemize}

\subsubsection*{Monday. Methods: trees and forests.}

\begin{itemize}
\item \bibentry{wager2018estimation}\vskip 2 ex
\item \textit{Optional/reference:} \bibentry{athey2016recursive}
\item \textit{Optional/reference:} \bibentry{athey2019generalized}\vskip 2 ex
\item \textit{Textbook reference:} \bibentry{hastie2009elements} Trees 9.2; Random forests, 15.1-15.3.
\item \textit{Textbook reference:} \bibentry{hernan2010causal} Treatment effect heterogeneity, 4.1-4.4.
\end{itemize}

\subsubsection*{Wednesday.}
\begin{itemize}
\item \bibentry{kunzel2019metalearners}\vskip 2 ex
\item \textit{Textbook reference:} \bibentry{hernan2010causal} Why model, 11.1-11.5. 
\end{itemize}


\subsection*{Week 4. Representation learning for causal inference.}

\subsubsection*{Monday.}
\begin{itemize}
\item \bibentry{johansson2016learning}
\end{itemize}

\subsubsection*{Wednesday.}
\begin{itemize}
\item \bibentry{shalit2017estimating}
\item \textit{Optional/reference:} \bibentry{nie2021quasioracle}\vskip 2 ex
\item \textit{Textbook reference:} \bibentry{hastie2009elements} Neural networks, 11.1-11.4.
\item \textit{Textbook reference:} \bibentry{bishop2006pattern} Neural networks, 5.\vskip 2 ex
\item \textit{Textbook reference:} \bibentry{hastie2009elements} Graphical representation of causal effects. 
\end{itemize}

\textcolor{cyan}{Proposed tutorial topics must be emailed to the instructor by EOD the Wednesday of Week 4.}

\subsection*{Week 5. When ignorability fails: proxies and instruments.}
\subsubsection*{Monday.}
\begin{itemize}
\item \bibentry{tchetgen2024proximal}
\item \textit{Optional/reference:} \bibentry{miao2018identifying} \vskip 2 ex
\item \textit{Textbook reference:} \bibentry{hernan2010causal} G-computation, 13, 14. 
\end{itemize}




\subsubsection*{Wednesday.}
\begin{itemize}
\item \bibentry{pmlr-v70-hartford17a}\vskip 2 ex
\item \textit{Textbook reference:} \bibentry{hastie2009elements} Instrumental variables, 16. 
\end{itemize}

\subsection*{Week 6. Policy learning and evaluation.}
\subsubsection*{Monday.}
\begin{itemize}
\item \bibentry{athey2021policy}
\item \textit{Optional/reference:}
\bibentry{kitagawa2018should}
\end{itemize}


\subsubsection*{Wednesday.}
\begin{itemize}
\item \bibentry{dudik2014doubly}
\end{itemize}

\subsection*{Week 7. Efficient estimation: DR, DML, and TMLE.}

\subsubsection*{Monday.}
\begin{itemize}
\item \bibentry{kennedy2022semiparametric}. \\
\textbf{Parts 1-3.}\vskip 2 ex
\item  \textit{Optional/reference:} Edward Kennedy's slides from the ACIC conference in 2022: \emph{Machine learning \& nonparametric efficiency
in causal inference} \url{https://www.ehkennedy.com/uploads/5/8/4/5/58450265/tutorial.pdf}\\
The slides present the same material as the paper, with some additional narrative \& figures, and references throughout to relevant historical work. There are also some videos of Kennedy giving lectures on this topic on YouTube, e.g., \url{https://youtu.be/vZKdObmilQU}.
\item  \textit{Textbook reference:} \bibentry{hampel1974influence}
\item  \textit{Textbook reference:} \bibentry{lehmann1999elements}\\
\textbf{In particular, sections 1.1-1.4 on limits and rates of convergence. Fisher Information is covered in sections 7.1 and 7.2. The theory of statistical functionals and the influence curve is covered in sections 6.2 and 6.3.}
\item \textit{Textbook reference:} \bibentry{hastie2009elements} Doubly robust machine learning estimators, 18.4. 
\end{itemize}

\subsubsection*{Wednesday.}
\begin{itemize}
\item \bibentry{chernozhukov2018double} \\
\textbf{Core sections/excerpts to be assigned.}
\end{itemize}

\textcolor{cyan}{Meet with the instructor to discuss work-in-progress by the end of Week 7.}


\subsection*{Week 8. Generalizability, transportability, and policy shift.}

\subsubsection*{Monday.}
\begin{itemize}
\item \bibentry{degtiar2023review}
\item \textit{Optional/reference:} \bibentry{pearl2014external}\vskip 2 ex
\item \textit{Textbook reference:} \bibentry{hastie2009elements} Discussion of transportability 4.3. 
\end{itemize}

\subsubsection*{Wednesday.}
\begin{itemize}
\item \bibentry{athey2021matrix}
\end{itemize}

\subsection*{Week 9. AI-assisted measurement and valid inference.}

\subsubsection*{Monday.}
\begin{itemize}
\item \bibentry{zrnic2024cross}
\item \textit{Optional/reference:} \bibentry{angelopoulos2023prediction}
\end{itemize}

\subsubsection*{Wednesday.}
\begin{itemize}
\item \bibentry{egami2024using}
\item \textit{Optional/reference:} \bibentry{halterman2025codebook}
\end{itemize}

\textcolor{cyan}{Final materials due by 11:59 PM on the Friday of finals week.}

\bibliography{\plscbibpath}
\end{document}
